\section{Descripción del problema}

\subsection{Problemática general}
En programas universitarios con mallas curriculares complejas, el estudiantado se enfrenta a la tarea de planear su ruta de estudio en un entorno lleno de restricciones de prerrequisitos, límites de créditos por periodo y oferta variable de asignaturas. El \textbf{gran problema} que se busca resolver es la ausencia de una herramienta inteligente que, a partir de la estructura del plan de estudios y del historial académico de cada estudiante, genere rutas académicas coherentes y eficientes que reduzcan el riesgo de retrasos y prolongación innecesaria del tiempo de graduación.

Resolver este problema es importante porque las decisiones de inscripción de asignaturas tienen un impacto directo en el tiempo total de formación, la carga académica asumida en cada semestre y la probabilidad de enfrentar cuellos de botella derivados de prerrequisitos mal secuenciados. Sin apoyo sistemático, los estudiantes suelen depender de recomendaciones informales o de una exploración manual del currículo, lo que aumenta la probabilidad de rutas ineficientes, pérdida de motivación y posibles retrasos en la culminación del programa. Un planificador inteligente de ruta de estudio contribuye a mitigar estos riesgos al ofrecer planes consistentes con las reglas curriculares, acordes con las preferencias del estudiante y orientados a optimizar su trayectoria académica.

\pagebreak